\documentclass[11pt]{article}

\usepackage[a4paper,margin=1in]{geometry}
\usepackage{microtype}
\usepackage{amsmath,amssymb,amsthm}
\usepackage[hidelinks]{hyperref}
\usepackage{booktabs}
\usepackage{enumitem}

\title{RQ-VAE and DRQ: Coupled versus Decoupled Residual Quantization}
\author{}
\date{}

\begin{document}
\maketitle

\section{The point}

RQ-VAE trains two things in one loop---a continuous geometry (the encoder and
decoder) and a discrete code assignment (the codebook)---and ties them together
with the straight-through estimator. DRQ does the same two jobs in two separate
stages: first learn the geometry with a plain continuous VAE, then assign
discrete codes with hierarchical $k$-means on the frozen geometry. The split
matters because the two jobs pull in opposite directions. Learning geometry
wants the latent vectors spread over the whole manifold; assigning codes wants
an even, flat distribution over codes. One coupled objective cannot do both at
once, so RQ-VAE ends up with far fewer usable codes than it nominally has, and
items share the same code. DRQ's two stages each satisfy their own demand.

\section{RQ-VAE}

An encoder maps an item embedding $x$ to a continuous latent vector
\begin{equation}
  z = E(x),
  \label{eq:enc}
\end{equation}
and a decoder $D$ maps a quantized copy of $z$ back to a reconstruction
$\hat x$. Quantization is residual and hierarchical, with $L$ codebooks
$\mathcal{C}^{(0)}, \dots, \mathcal{C}^{(L-1)}$, each holding $K$ codewords.
Starting from $r_0 = z$, level $\ell$ assigns the nearest codeword and subtracts
it:
\begin{align}
  c^{(\ell)} &= \operatorname*{arg\,min}_{c \in \mathcal{C}^{(\ell)}}
                \bigl\| r_{\ell-1} - c \bigr\|_2^2,
                \label{eq:assign} \\
  r_\ell &= r_{\ell-1} - c^{(\ell)}.
  \label{eq:residual}
\end{align}
After $L$ levels the reconstruction and the loss are
\begin{align}
  \hat z &= \sum_{\ell=0}^{L-1} c^{(\ell)},
  \qquad \hat x = D(\hat z),
  \label{eq:recon} \\
  \mathcal{L} &= \| x - \hat x \|_2^2
  + \sum_{\ell=0}^{L-1} \Bigl(
      \bigl\| \mathrm{sg}(r_{\ell-1}) - c^{(\ell)} \bigr\|_2^2
    + \beta \bigl\| r_{\ell-1} - \mathrm{sg}(c^{(\ell)}) \bigr\|_2^2
    \Bigr).
  \label{eq:loss}
\end{align}
Here $\mathrm{sg}$ is stop-gradient and $\beta$ is the commitment weight. The
assignment \eqref{eq:assign} is not differentiable. The straight-through
estimator therefore treats the quantized vector as the identity: the decoder's
gradient with respect to $\hat z$ is copied to $z$, so the encoder receives
$\partial \mathcal{L}/\partial \hat z$ unchanged. Implemented as
$z_{\mathrm{ste}} = z + \mathrm{sg}(\hat z - z)$, and the decoder runs on
$z_{\mathrm{ste}}$.

The item's semantic ID is the tuple $(c^{(0)}, \dots, c^{(L-1)})$, with nominal
capacity $K^L$. The set of all possible reconstructions $\hat z$ is the Minkowski
sum $\mathcal{C}^{(0)} + \cdots + \mathcal{C}^{(L-1)}$: a grid of $K^L$ points,
each the sum of one codeword per level.

\section{Why collisions happen}

Nominal capacity overstates what the codebook can actually separate. Two
different items collide when they receive the same tuple, and the collision rate
is set by how many tuples are in real use, not by $K^L$.

DRQ quantifies this with an expected overlap rate. Let $\pi_i$ be the fraction
of items assigned to codeword $i$, and $\sigma^2$ the variance of a small
perturbation applied to the latent vectors at retrieval time. The overlap rate
and the effective codebook size are
\begin{equation}
  O_\pi = \sum_{i} \sum_{j} \pi_i \pi_j
          \exp\!\Bigl( -\frac{\| c_i - c_j \|_2^2}{4 \sigma^2} \Bigr),
  \qquad
  K_{\mathrm{eff}} = \frac{1}{O_\pi}.
  \label{eq:overlap}
\end{equation}
Splitting \eqref{eq:overlap} into diagonal and off-diagonal terms shows the two
failure modes:
\begin{equation}
  O_\pi =
    \underbrace{\sum_{i} \pi_i^2}_{\text{distribution floor}}
    + \underbrace{\sum_{i \ne j} \pi_i \pi_j
      \exp\!\Bigl( -\frac{\| c_i - c_j \|_2^2}{4 \sigma^2} \Bigr)}_{\text{geometry cross-overlap}}.
  \label{eq:split}
\end{equation}
For a uniform codebook with infinitely separated codewords, $\pi_i = 1/K$ and
the off-diagonal term vanishes, leaving $O_\pi = 1/K$ and $K_{\mathrm{eff}} = K$.
Any deviation from that ideal shrinks $K_{\mathrm{eff}}$:

\begin{itemize}[nosep]
  \item \textbf{Distribution penalty.} When usage concentrates on a few codewords,
    $\sum_i \pi_i^2$ rises above $1/K$. This floor remains even when the
    codewords sit far apart.
  \item \textbf{Geometry penalty.} The Minkowski-sum grid is flat. When the data
    lie on a curved manifold, most grid points sit where no item lives, and the
    codewords that do carry items stay close together, keeping the cross-overlap
    large.
\end{itemize}

Both effects push $K_{\mathrm{eff}}$ far below $K^L$. Then $4.8$M items can
collide even though $1024^3 \approx 1.07 \times 10^9$ looks like plenty of room.

\section{DRQ}

DRQ keeps the hierarchical code structure but runs the two jobs separately.

\textbf{Stage 1: continuous geometry.} Train the encoder and decoder as a plain
VAE, with no codebook, no straight-through estimator, and no commitment term:
\begin{equation}
  \mathcal{L} = \mathbb{E}_{x} \bigl\| x - D(E(x)) \bigr\|_2^2
              + \lambda \mathcal{L}_{\text{contrastive}}.
  \label{eq:drq-stage1}
\end{equation}
The contrastive term is InfoNCE over positive pairs $(a, b)$ of related items
(co-viewed or semantically similar), with in-batch negatives and cosine
similarity scaled by a temperature $\tau$:
\begin{equation}
  \mathcal{L}_{\text{contrastive}} = -\frac{1}{B}\sum_{k=1}^{B}
    \log\frac{\exp(\mathrm{sim}(z_{a_k}, z_{b_k})/\tau)}
             {\sum_{j=1}^{B}\exp(\mathrm{sim}(z_{a_k}, z_{b_j})/\tau)},
  \qquad z = E(x).
  \label{eq:drq-infonce}
\end{equation}
The paper's DRQ-VAE sets $\lambda = 0$, so stage 1 is reconstruction only. The
DRQ-VAE+CL variant uses $\lambda = 0.3$. The contrastive term is symmetric in
practice and $\tau$ is not stated in the version we checked, so our default
$\tau = 0.1$ is an assumption.
Without a grid forcing it, the latent space is free to spread along the manifold.

\textbf{Stage 2: discrete assignment.} Freeze the geometry and run hierarchical
$k$-means on the latent vectors. Level $0$ clusters $z$ itself; each later level
clusters the residual left by the level above. Each level's centroids are fitted
on a uniform random sample $\mathcal{S}$ of latents (default $10^6$ rows), then
every latent in the corpus is assigned to its nearest centroid, as in
\eqref{eq:assign}. Centroids are means, not gradient-updated:
\begin{equation}
  c^{(\ell)}_i = \frac{1}{|\mathcal{S}^{(\ell)}_i|}
                 \sum_{r \in \mathcal{S}^{(\ell)}_i} r,
  \qquad
  \mathcal{S}^{(\ell)}_i = \bigl\{ r \in \mathcal{S} : r \text{ assigned to centroid } i \bigr\}.
  \label{eq:drq-mean}
\end{equation}
Within each Lloyd iteration every centroid is updated from all of its sampled
members, so rarely used codes still get updates, unlike codes trained only
through the straight-through estimator. This is an argument, not a measured
result. Empty clusters are not special-cased in our formulation.

\section{What changed}

\begin{itemize}[nosep]
  \item RQ-VAE optimizes geometry and assignment together; DRQ optimizes geometry
    first and assignment second.
  \item The overlap kernel width $\sigma$ is a free choice in our implementation.
    We set $\sigma = 0.1 \times$ the median nearest-neighbour distance between
    used centroids, per level. The paper's $\sigma$ is defined as a
    retrieval-time perturbation variance, and our choice is not derived from it.
  \item RQ-VAE's assignment gradients flow only through selected codes via the
    straight-through estimator; DRQ's assignment is a full-dataset clustering.
  \item RQ-VAE forces the latent space onto a Minkowski-sum grid; DRQ lets the
    latent space spread first, then adapts the grid to where the data are.
\end{itemize}

On the paper's industrial dataset (over $15$M items, $K = 4096$, $L = 3$),
RQ-VAE used only $512$ of $4096$ codes at the coarsest level, and DRQ cut the
collision rate from $0.229$ to $0.075$. The collision-rate denominator is not
stated in the version we checked, so treat these as reported values.

\section{Takeaways}

\begin{itemize}[nosep]
  \item Collisions are a capacity problem: the effective codebook size, not the
    nominal $K^L$, decides whether items collide.
  \item Two things shrink effective capacity: skewed code usage (distribution)
    and a flat grid on a curved manifold (geometry).
  \item Fixing either one inside a coupled objective is hard, because the two
    demand opposite behaviors.
  \item DRQ's answer is to decouple: learn geometry without codes, then assign
    codes without gradients.
\end{itemize}

\end{document}
